\section{The price of perfection}\label{sec:intro}
Which numbers can be the entanglement price of perfection for one fixed finite game? The exponent is the rate at which entanglement must grow as a fixed finite device is asked to play closer to perfection; if it can be uncomputable, no algorithm reading the finite rules of the game can predict that cost. Here prediction means computing the limiting exponent. For the entropy of a multidimensional shift of finite type the analogous question has a computability-theoretic answer~\cite{HM10}; here the quantity is the exponent of the least local dimension, and of the usable entanglement, needed to lose at most $\eps$. Every positive left-computably-enumerable real occurs, including uncomputable reals. The exponents also go strictly beyond the left-c.e. reals: $(4-\theta_K)/4$ is a noncomputable right-c.e. example, where $\theta_K=\sum_{e\in K}4^{-(e+1)}$ and $K$ is the halting set. The same games have their dimension exponent as their logarithmic usable-entanglement coefficient. Their local-reset heat has twice that coefficient, while their joint-reset heat stays bounded.

A real is \emph{left-computably-enumerable}, or \emph{left-c.e.}, if it is the limit of a computable nondecreasing sequence of rationals. Increasing approximations below are understood in this non-strict sense. All logarithms without a displayed base are base two, and entanglement is measured in ebits. For a finite game $G$ of quantum value one, let $D_G(\eps)$ be the least local dimension of a strategy with loss at most $\eps$, padding the two local spaces equally. Define
\begin{equation}\label{eq:usable-definition}
 d^\to_G(\eps)=\inf_{\text{finite-dimensional strategies of loss }\le\eps}
 E_D^{A\to B}(\rho_{AB}).
\end{equation}
The state may be mixed. Its one-way distillable entanglement is the asymptotic Bell-pair rate using local operations and classical communication from Alice to Bob. A real is \emph{right-c.e.} if its negative is left-c.e.; it is $\Delta^0_2$ if computable relative to the halting set.

\begin{theorem}[Realization]\label{thm:realization}
For every positive left-c.e. real $\alpha$ there is a fixed finite two-player Boolean constraint-system game $G_\alpha$, with rational question probabilities, quantum value one and no finite-dimensional perfect strategy, such that, as $\eps\downarrow0$,
\begin{align}
 \log_2D_{G_\alpha}(\eps)&=\alpha\log_2(1/\eps)+o(\log(1/\eps)),\label{eq:realization-dimension}\\
 d^\to_{G_\alpha}(\eps)&=\alpha\log_2(1/\eps)+o(\log(1/\eps)).\label{eq:realization-usable}
\end{align}
There is also such a game $\widehat G_K$ satisfying both laws with $\alpha=(4-\theta_K)/4$, which is right-c.e. and not left-c.e. Thus the positive dimension exponents in this BCS class strictly contain the positive left-c.e. reals and are contained in the positive $\Delta^0_2$ reals. Every finite dimension exponent of a rational finite game is $\Delta^0_2$ whenever its defining limit exists.
\end{theorem}

A finite game's price of precision can have any positive left-c.e. coefficient and can also be non-left-c.e. The coefficient measures usable entanglement in the original mixed state. For the halting examples, local-reset heat has an uncomputable logarithmic coefficient, while joint-reset heat stays bounded. The precise heat statements are Corollary~\ref{cor:heat} and Theorem~\ref{thm:gate-resources}.

\subsection{The flagship halting game}
Fix a standard effective enumeration of programs, with halting set $K\subseteq\N=\{0,1,\ldots\}$.
\begin{theorem}[The version 1.0.0 game]\label{thm:main}
There is a fixed finite two-player Boolean constraint-system game $G_K$, with rational question probabilities, such that $\omega_q(G_K)=1$, no finite-dimensional strategy wins perfectly, and
\begin{equation}\label{eq:main}
 D_{G_K}(\eps)=\eps^{-(3+\thetaK)/4+o(1)},\qquad
 \thetaK=\sum_{e\in K}4^{-(e+1)}.
\end{equation}
The exponent is Turing-equivalent to $K$. The sequence $\log D_{G_K}(2^{-n})/n$ has no computable certified rate of convergence.
\end{theorem}

This is the original explicit degree example. The existence of increasing computable rational sequences with uncomputable limits is classical, beginning with Specker~\cite{Spe49}; the construction places such a limit inside the resource law of a single finite game. Sections~\ref{sec:prelim}--\ref{sec:game} retain the construction and proof of~\cite{DM-v1}. Sections~\ref{sec:realization}--\ref{sec:heat} extend it. Section~\ref{sec:limits} proves the exact-trace barrier; Section~\ref{sec:gate} escapes it and states the remaining classification questions.

\subsection{Related work and credit}
Mousavi, Nezhadi and Yuen prove that deciding whether the quantum value is exactly one is $\Pi^0_2$-complete~\cite{MNY20,MNY22}. In particular, there is no total computable dimension bound, uniform over all value-one games and all precisions. Indeed, such a bound would make value one a universal integer quantification of decidable bounded-dimensional feasibility statements, placing that problem in $\Pi^0_1$. Their obstruction is uniform over games. Theorems~\ref{thm:realization} and~\ref{thm:main} instead identify convergent coefficients for individual fixed games. The theorem $\mathrm{MIP}^*=\mathrm{RE}$ puts halting information in the value of a computable family of games~\cite{JNVWY}; it does not by itself identify a resource exponent of a fixed value-one game.

The algebraic connection between constraint games and operator assignments goes back to Cleve and Mittal~\cite{CM14} and, for linear systems, Cleve, Liu and Slofstra~\cite{CLS17}. Slofstra established an embedding theorem for solution groups~\cite{Slo20} and the non-closure of quantum correlations~\cite{Slo19}. Slofstra and Vidick introduced the hyperlinear profile as a quantitative bridge from approximate representations to dimension requirements~\cite{SV18}. Slofstra later obtained a finitely presented group with at least subexponential hyperlinear profile~\cite{Slo18}. Ji, Leung and Vidick construct a fixed finite three-player game of value one for which near-perfect play requires at least $\Omega(\eps^{-c})$ qubits for a fixed $c>0$~\cite{JLV20}. Coladangelo constructs a fixed finite two-player score game for which every $\eps$-optimal strategy has local dimension $2^{\Omega(\eps^{-1/8})}$~\cite{Col20}. These results supply the framework and important growth examples; they do not identify the matched finite exponent in~\eqref{eq:main}.

Mehta, Slofstra and Zhao construct halting-dependent Pauli systems and embed them in finite presentations with polynomial bounds on relation decompositions~\cite[Proposition~4.8]{MSZ23}. Those bounds transfer lower estimates; they do not assert small simultaneous finite-model extensions. The construction here supplies its own upper models. The conjugacy gadgets used in our compiler are those of Fanizza, Kroell, Mehta, Paddock, Rochette, Slofstra and Zhao~\cite[Lemma~5.8 and Theorem~5.5]{FKMPRSZ25}. Their credit is essential: the passage from nested conjugacy to ordinary Boolean contexts is not a new gadget of this paper. We give its matrices and converse calculation because quantitative completeness is needed here.

The state-preserving Boolean correspondence was established in \emph{The price of precision}~\cite[Theorem~4.9]{DM-precision}; its linear-system predecessor and further dimension examples appear in the companion paper~\cite{DM-game}. We use that correspondence as an existing theorem. Section~\ref{sec:game} records the elementary square-root estimate needed for the present local-dimension application and checks the constants' effect on the exponent. Theorem~\ref{thm:realization} proves and strengthens the computable-coefficient inclusion conjectured in the discussion of~\cite[Problem~18.3(i)]{DM-precision}, for its original usable-entanglement demand as well as dimension. The $\Delta^0_2$ upper bound for dimension exponents is~\cite[Proposition~11.13]{DM-precision}; Lemma~\ref{lem:computability} records the rational case.

Uncomputable asymptotic data in finite systems have other precedents. Hochman and Meyerovitch characterize the entropies of multidimensional shifts of finite type as exactly the nonnegative right-c.e. reals~\cite{HM10}; their construction passes through a shift whose maximal symbol frequency is a prescribed right-c.e. real. Their theorem is the closest template for identifying the possible values of a physical quantity of a fixed finite system with a computability class. Here every positive left-c.e. real is realized, together with a non-left-c.e. right-c.e. family; equality for the full game spectrum is not asserted. Sapir, Birget and Rips relate group isoperimetry and computation~\cite{SBR02}; a Dehn exponent is not a dimension exponent, and no identification between them is used here. Cubitt, Perez-Garcia and Wolf prove undecidability of the spectral gap over families of local Hamiltonians~\cite{CPGW22}. Bausch, Cubitt and Watson obtain uncomputable phase diagrams for a continuous one-parameter family~\cite{BCW21}. These are precedents for asymptotic noncomputability, rather than proofs of~\eqref{eq:main}.

The binary adding machine used below is the classical odometer on a Cantor set; its interpretation through operator algebras is standard~\cite{HSWZ13}. The new point needed here is a finite polynomial presentation and a carry estimate that remain quantitative in a non-flat state. The final information bound uses Holevo's theorem~\cite{Hol73} and the random-access-code principle exemplified by Nayak~\cite{Nay99}. Its branching estimate is proved explicitly.


The density classifications of Jockusch--Schupp and Downey--Jockusch--Schupp~\cite{JS12,DJS13} are prior work on computation encoded in densities. The moving-window estimate in Section~\ref{sec:realization} uses the binary radical-inverse ordering of van der Corput~\cite{vdC35}, not a new discrepancy principle. Vidick's \emph{Almost synchronous quantum correlations}~\cite{Vid22} supplies the closely related spectral-compression mechanism; Section~\ref{sec:limits} gives the word-moment form needed here. Tracial moment optimization~\cite{BK12,KMV22} supplies the exact-trace lower hierarchy. Its distinction from general commuting-model NPA~\cite{NPA08} is the distinction used to escape the barrier in Section~\ref{sec:gate}.

Devetak--Winter hashing~\cite{DW05} is the operational last step in Section~\ref{sec:usable}. The new ingredient is the inactive-outcome extension of the triangular extraction in~\cite[Theorem~6.2]{DM-precision} and its cyclic finite-alphabet form~\cite[Lemma~6.3]{DM-precision}, with random targets and an original-mixed-state bridge. The heat conclusion invokes the existing reset sandwich of~\cite[Theorem~2.1 and Corollary~2.2]{DM-heat}. Kolchinsky and Wolpert~\cite{KW20} study noncomputable heat functions for universal Turing machines. Corollary~\ref{cor:heat} and Theorem~\ref{thm:gate-resources} identify uncomputable logarithmic precision coefficients for fixed finite nonlocal games and compare their local and joint reset schedules. The halting-degree dimension construction, its compiler, clocks and proofs appeared in version 1.0.0~\cite{DM-v1}; they are retained here. The new realization, extraction and consequences are proved in the added sections. Section~\ref{sec:gate} additionally uses the finite perfect-completeness form of Ji--Natarajan--Vidick--Wright--Yuen, states its exact theorem numbers, and credits the Connes embedding and commuting-complexity context, including Lin~\cite{Lin25}.

\subsection{Proof outline}
At source precision $4^{-N}$, three compulsory clocks cost $3N-o(N)$ bits of logarithmic dimension. They hold an address, fuel and time. For $G_K$, a finite counter program recognizes halting inputs on separated base-four cylinders. Reset forces initialized computations at every finite fuel scale. Exact tracial models must flag every halting cylinder; compactness extends this to all weighted approximate models, without an effective modulus.

A conditional clock contributes $\theta_KN-o(N)$ further bits. Explicit finite histories are diagonal labels on the existing registers. A moving register of optional bits realizes the conditional clock in dimension $2^{\theta_NN+O(\log N)}$, with computable $\theta_N\uparrow\theta_K$. Source precision $4^{-N}$ becomes game loss of order $16^{-N}$, giving the denominator four.

For an arbitrary left-c.e. $\theta\in[0,1]$, an address reader accepts exactly the binary values below $\theta$. Radical-inverse window counts match its finite flag density. Extra compulsory clocks add to the numerator, and replacing polynomial residuals by fixed powers divides the exponent. These two operations cover every positive left-c.e. $\alpha$. Finally a coherent extraction with an inactive outcome converts the expected active count into distillable Bell pairs on the original mixed state. The cited heat sandwich then matches both reset laws.


The complementary construction reverses the flag rule: an unflagged decoded cylinder must carry a finite perfect witness for its halting game. A fixed counter transducer with two operator registers tests every relation of that witness's projection presentation, packed into three data generators. Nonhalting unflagged corners have zero weight in matrix-approximable traces, and moment compression transfers that bound to the actual strategy state. Finite perfect witnesses can be admitted so slowly that their combined logarithmic dimension is at most $\sqrt N$, while the unflagged density tends to $\theta_K$. The flag density is therefore $1-\theta_K$, giving exponent $(4-\theta_K)/4$. The same extraction and heat sandwich apply. The exact-trace minimum is strictly smaller, as the barrier forces it to be.
